\begin{problem}[25 pts]
\statementfigure{%
It is one-dimensional motion along x-axis. Figure on the right shows
change in an acceleration with time. Find the mean velocity averaged
over time interval $[\qty{2}{\second},\qty{10}{\second}]$, if the initial
velocity equals to $\qty{-8}{\metre\per\second}$?
}{\resizebox{\linewidth}{!}{\begin{tikzpicture}
\begin{axis}[width=10.5cm,height=4.5cm,axis lines=middle,
 xmin=0,xmax=11,ymin=-0.8,ymax=9.5,
 xlabel={$t$ (\unit{s})},ylabel={$a$ (\unit{\metre\per\second\squared})},
 xtick={0,2,4,6,8,10},ytick={0,2,4,6,8},
 tick label style={font=\large},label style={font=\large},
 ylabel style={at={(axis description cs:0,1)},anchor=south west,
 rotate=0,yshift=3pt},
 grid=major,grid style={PhysicsLine},clip=false]
\addplot[physics line] coordinates
 {(0,0)(2,0)(2,4)(4,4)(4,0)(8,0)(10,8)};
\draw[dashed,PhysicsBlue] (axis cs:10,0)--(axis cs:10,8)--(axis cs:0,8);
\end{axis}
\end{tikzpicture}
}}
\end{problem}
\begin{solution}
\qstep{1}{Name the intervals and accelerations.}
The intervals 2--4, 4--8 and 8--10 s have durations
$\Delta t,2\Delta t,\Delta t$.
Let $a_1$ be the acceleration on 2--4 s and $a_{\rm f}$ the acceleration at the end of the ramp.
Before 2 s, acceleration is zero, so the velocity is still its initial value $v_0$.
After the first interval it is
$v_1=v_0+a_1\Delta t$. It then stays constant until 8 s.
\qstep{2}{Find the displacement in each interval.}
During 8--10 s, let $t'$ be time since 8 s, with $0\leq t'\leq\Delta t$.
Acceleration rises linearly from zero:
\[
a(t')=a_{\rm f}\frac{t'}{\Delta t},\qquad
v(t')=v_1+\frac{a_{\rm f}(t')^2}{2\Delta t}.
\]
Integrating velocity gives the three signed displacements:
\[
\begin{aligned}
\Delta x_1&=\frac{v_0+v_1}{2}\Delta t,\qquad
\Delta x_2=2v_1\Delta t,\\
\Delta x_3&=\int_0^{\Delta t}v(t')\,dt'
=v_1\Delta t+\frac{a_{\rm f}(\Delta t)^2}{6}.
\end{aligned}
\]
\qstep{3}{Divide total displacement by total time.}
\[
\bar v=\frac{\Delta x_1+\Delta x_2+\Delta x_3}{4\Delta t}
=\frac{\frac{v_0+v_1}{2}\Delta t+3v_1\Delta t+
\frac{a_{\rm f}(\Delta t)^2}{6}}{4\Delta t}.
\]
\qstep{4}{Calculate and check the signed areas.}
Use $\Delta t=2$, $v_0=-8$, $a_1=4$ and $a_{\rm f}=8$ in SI units.
Then $v_1=-8+4(2)=0$, giving displacements
$-8$, $0$ and $\frac{8(2)^2}{6}=\frac{16}{3}\,\unit{\metre}$.
\[
\bar v=\frac{-8+0+\frac{16}{3}}{8}\,\unit{\metre\per\second}
=-\frac{1}{3}\,\unit{\metre\per\second}
\approx-\qty{0.333}{\metre\per\second}.
\]
\sourcefig{0.59\linewidth}{velocity-v3}
The final branch is curved, so a trapezoid would give the wrong area.
Signed displacements partly cancel; mean speed is a different quantity.
\end{solution}

\ifphysicsshowsolutions\clearpage\fi
\begin{solution}
\subsection*{Problem 1: alternative using the original time}
\qstep{1}{Write the acceleration on 8--10 s.}
The straight line joins $(8,0)$ and $(10,8)$ on the acceleration graph.
Using $t$ in seconds and numerical values in SI units,
\[
a(t)=\frac{8-0}{10-8}(t-8)=4(t-8),\qquad 8\leq t\leq10.
\]
The coefficient $4$ has units \unit{\metre\per\second\cubed},
so $a$ is in \unit{\metre\per\second\squared}.
\qstep{2}{Integrate acceleration to find velocity.}
On 2--4 s, velocity rises from $-\qty{8}{\metre\per\second}$ to
$-8+4(4-2)=0$; it then stays zero until 8 s. Thus $v(8)=0$ and
\[
\begin{aligned}
v(t)&=v(8)+\int_8^t a(t')\,dt'\\
    &=\int_8^t4(t'-8)\,dt'=2(t-8)^2,\qquad 8\leq t\leq10.
\end{aligned}
\]
Here $t'$ is the integration variable; $v$ is in \unit{\metre\per\second}.
\qstep{3}{Integrate velocity to find displacement.}
\[
\Delta x_{8\text{--}10}=\int_8^{10}v(t)\,dt
=\left[\frac{2(t-8)^3}{3}\right]_8^{10}
=\frac{16}{3}\,\unit{\metre}.
\]
\qstep{4}{Include the earlier intervals and find the mean.}
The signed area on 2--4 s is
$\Delta x_{2\text{--}4}=\frac{-8+0}{2}(4-2)=-\qty{8}{\metre}$.
On 4--8 s, $\Delta x_{4\text{--}8}=0$ because the object is at rest.
Therefore,
\[
\bar v=\frac{\Delta x_{2\text{--}4}+\Delta x_{4\text{--}8}
+\Delta x_{8\text{--}10}}{10-2}
=\frac{-8+0+\frac{16}{3}}{8}\,\unit{\metre\per\second}
=-\frac{1}{3}\,\unit{\metre\per\second}.
\]
\textbf{Answer:} $\bar v\approx-\qty{0.333}{\metre\per\second}$,
the same as in the preceding method.
\end{solution}

\quiznextproblem
\begin{problem}[25 pts]
The tram starts moving with the constant acceleration of
\qty{0.3}{\metre\per\second\squared}. Find the minimum passenger's velocity,
when the passenger could wait for \qty{8}{\second} and then catch up the
tram by running after.
\end{problem}
\begin{solution}
\qstep{1}{Choose the model and sketch the minimum-speed case.}
Assume a common starting position, a tram starting from rest, and a
passenger who waits there for time $t_{\rm w}$ and then runs at constant speed $v_{\rm P}$.
The source does not explicitly specify constant running speed.
\sourcefig{0.48\linewidth}{tram-limit}
At the minimum speed the passenger's position line touches the tram's parabola.
\qstep{2}{Write the meeting equation.}
For $t>t_{\rm w}>0$,
\[
x_T=\frac{at^2}{2},\qquad x_P=v_{\rm P}(t-t_{\rm w}),\qquad
at^2-2v_{\rm P}t+2v_{\rm P}t_{\rm w}=0.
\]
\qstep{3}{Require a real meeting after the wait.}
The discriminant is nonnegative when $v_{\rm P}(v_{\rm P}-2at_{\rm w})\geq0$.
For positive running speed this gives
\[
v_{\rm P,min}=2at_{\rm w},\qquad v_{\rm P}\geq2at_{\rm w}.
\]
The formal boundary $v_{\rm P}=0$ is invalid: its only meeting time is $t=0<t_{\rm w}$.
At the valid boundary the double root is $t_*=\frac{v_{\rm P}}{a}=2t_{\rm w}$.
\qstep{4}{Calculate the speed.}
\[
v_{\rm P,min}=2(\qty{0.3}{\metre\per\second\squared})(\qty{8}{\second})
=\qty{4.80}{\metre\per\second}.
\]
The meeting occurs after 16 s in total, including 8 s of running.
This is the running speed, not an average including the wait.

\textit{Alternative:} minimize
$v_{\rm P}(t)=\frac{at^2}{2(t-t_{\rm w})}$ for $t>t_{\rm w}$.
Since $v_{\rm P}'(t)=\frac{at(t-2t_{\rm w})}{2(t-t_{\rm w})^2}$, the minimum is at
$t=2t_{\rm w}$, again giving $v_{\rm P,min}=2at_{\rm w}$.
\end{solution}

\quizpairbreak
\begin{problem}[25 pts]
\statementfigure{%
Two balls were thrown from the ground at angles of \ang{30} and
\ang{45} to the horizon. The ranges $R$ of the balls were the same.
Find the ratio of their maximum heights during the flights.
Neglect air resistance.
}{\resizebox{\linewidth}{!}{\begin{tikzpicture}[font=\small]
\draw[physics line,-{Latex}] (-.15,0)--(5.6,0);
\draw[physics line,dashed,domain=0:5,samples=40]
 plot (\x,{\x*(1-\x/5)});
\draw[physics accent,dashed,domain=0:5,samples=40]
 plot (\x,{tan(30)*\x*(1-\x/5)});
\draw (0.65,0) arc[start angle=0,end angle=30,radius=.65];
\draw (1.05,0) arc[start angle=0,end angle=45,radius=1.05];
\node[fill=white,inner sep=1pt] at (.8,.18) {$30^\circ$};
\node[fill=white,inner sep=1pt] at (.8,.62) {$45^\circ$};
\node[below] at (0,0) {$0$};
\node[below] at (5,0) {$R$};
\node at (2.8,-.7) {Schematic, not to scale};
\end{tikzpicture}
}}
\end{problem}
\begin{solution}
\qstep{1}{Identify the common range and the two heights.}
Use the equal launch and landing levels shown, with uniform downward
gravity. The initial speeds may differ.
\sourcefig{0.49\linewidth}{projectile-heights}
\qstep{2}{Write the height and range relations.}
At maximum height the vertical velocity is zero, so
\[
H_\theta=\frac{v_\theta^2\sin^2\theta}{2g},
\qquad R=\frac{v_\theta^2\sin(2\theta)}{g}.
\]
\qstep{3}{Eliminate the unknown launch speeds.}
Divide the relations:
\[
\frac{H_\theta}{R}
=\frac{\sin^2\theta}{2\sin(2\theta)}
=\frac{\tan\theta}{4}.
\]
Since both ranges are equal, the symbolic height ratio is
\[
\frac{H_A}{H_B}=\frac{\tan\theta_A}{\tan\theta_B}.
\]
\qstep{4}{Label the order and insert the angles.}
\[
\frac{H_{30}}{H_{45}}
=\frac{\tan30^\circ}{\tan45^\circ}
=\frac{1}{\sqrt3}\approx0.577.
\]
The labeled reverse ratio is equally valid:
$\frac{H_{45}}{H_{30}}=\sqrt3\approx1.73$.
Both are dimensionless; the \ang{45} trajectory is higher.

\textit{Alternative:} a parabola with roots $0,R$ and launch slope
$\tan\theta$ has equation
$y=\frac{\tan\theta}{R}x(R-x)$.
Its vertex at $x=\frac{R}{2}$ gives
$H_\theta=\frac{R\tan\theta}{4}$ and the same ratio.
\end{solution}

\quiznextproblem
\begin{problem}[25 pts]
A particle moves in uniform circular motion over a horizontal $xy$ plane.
At one instant, it moves through the point with coordinates
$[\qty{-10}{\metre},\qty{10}{\metre}]$ with velocity of
$(-20\hat{\mathbf i}+0\hat{\mathbf j})\,\unit{\metre\per\second}$
and acceleration of
$(0\hat{\mathbf i}+25\hat{\mathbf j})\,\unit{\metre\per\second\squared}$.
Find the position of the particle at moment, when its velocity equals to
$(20\hat{\mathbf i}+0\hat{\mathbf j})\,\unit{\metre\per\second}$.
\end{problem}
\begin{solution}
\qstep{1}{Locate the point with opposite velocity.}
\sourcefig{0.40\linewidth}{circle-v3}
Centripetal acceleration points upward toward the center.
The required velocity is opposite to the initial one, so the required
point is diametrically opposite. The motion is clockwise.
\qstep{2}{Find the radius from centripetal acceleration.}
Let $v_0=|\mathbf v_0|$ and $a_c=|\mathbf a_0|$.
For uniform circular motion,
\[
a_c=\frac{v_0^2}{R}\quad\Longrightarrow\quad R=\frac{v_0^2}{a_c}.
\]
\qstep{3}{Move upward by one diameter.}
The opposite point has the same $x$ coordinate and is $2R$ higher:
write the given initial coordinates as $(x_0,y_0)$. Then
\[
x_1=x_0,\qquad y_1=y_0+2R=y_0+\frac{2v_0^2}{a_c}.
\]
\qstep{4}{Calculate the radius and position.}
\[
R=\frac{20^2}{25}\,\unit{\metre}=\qty{16}{\metre},\qquad
x_1=-\qty{10}{\metre},\qquad
y_1=10+2(16)=\qty{42}{\metre}.
\]
\textbf{Answer:} $(x_1,y_1)=(-10,42)\,\unit{\metre}$.
Repeated revolutions give this same position, not additional coordinates.

\textit{Vector alternative:} the center is
$\mathbf r_C=\mathbf r_0+\frac{v_0^2}{a_c^2}\mathbf a_0$.
Opposite-point symmetry gives $\mathbf r_1=2\mathbf r_C-\mathbf r_0$,
which produces the same coordinates.
\end{solution}
